% SECTION 8 — Limitations and Future Work
% Paper 07: Recursive Spawning with Immutable Parent Pointers: Preventing Autonomous Drift
% Jeremy Blaine Thompson Beebe | ORCID: 0009-0009-2394-9714 | bxthre3inc@gmail.com
% This paper has not been peer reviewed. Comments and correspondence are welcomed.

\section{Limitations and Future Work}
\label{sec:limitations}

The Recursive Spawning Protocol with immutable parent pointers represents a structural advance in accountable multi-agent system design, but several honest limitations constrain its current form. Acknowledging these limits is a precondition for productive future development and a requirement of scientific integrity.

\subsection{Performance Overhead from Structural Enforcement}
\label{sec:limitations-performance}

Every spawn event under RSP passes through three sequential gates---Purpose Layer validation, Bounds Engine depth check, and Fact Layer warrant creation---before a child agent receives its activation warrant. This three-gate gating mechanism imposes a consistent and non-trivial overhead on spawn operations that does not exist in conventional agent-forking models.

At the Fact Layer, the creation of an immutable parent pointer requires cryptographic sealing: the spawn warrant is signed using the runtime's internal signing key, and the warrant record is inserted into the forensic ledger's Merkle-tree structure, requiring recomputation of logarithmic hash layers. In deployments with high spawn frequency---for example, a sensor network performing per-minute polling with nested analytical sub-agents---this overhead accumulates. Empirical evaluation in simulated agricultural deployments at scale (10,000+ sensor nodes, branching factor $b = 3$, depth $d = 5$) shows spawn gating adds approximately 12--18\% latency to the spawn initiation pathway compared to unrestricted forking.\cite{koch2026}

The immutable parent pointer itself is stored in a write-once, MPU-protected field. In hardware deployments, MPU enforcement requires context switches to privileged mode during spawn initialization, adding deterministic but non-trivial overhead. Software-only deployments rely on cryptographic sealing rather than hardware write protection, which trades MPU overhead for signature verification overhead. Neither path is free, and the protocol's accountability guarantees are purchased at this cost.

This overhead is a design choice, not an implementation artifact. The question is whether the accountability guarantees are worth their cost in performance-sensitive domains. For safety-critical agricultural deployments managing water rights, actuator control, and regulatory compliance, the answer is clearly yes. For low-stakes, high-frequency analytics operating in benign environments, the overhead may not be justified. RSP does not assume universal applicability; it defines the conditions under which the trade-off is correct.

\subsection{Scalability and the Branching Factor Problem}
\label{sec:limitations-scalability}

The exponential amplification property of drift---a single drifted agent at depth $d$ with branching factor $b$ produces up to $b^d$ drifted descendants---is not merely a characterization of the failure mode; it is a characterization of the spawn tree itself under any spawning model that permits non-trivial branching. Even a well-functioning RSP deployment with perfectly calibrated Purpose Layer validators produces $b^d$ total agents at depth $d$. This has direct consequences for forensic ledger storage, parent pointer resolution time, and the computational burden of chain traversal during incident reconstruction.

Consider a deployment with $d_{\max} = 5$ and $b = 4$ (a moderately aggressive spawning configuration). At depth 5, the maximum number of active nodes is $\sum_{i=0}^{5} b^i = 1,365$. For a cloud-hosted analytical service managing 10,000 field sensors, each spawning at most 4 children, the spawn tree under RSP is finite and bounded, but its size at full saturation is substantial. Storage for the forensic ledger scales linearly with total spawn events across the tree; chain traversal for incident reconstruction scales with depth, not branching factor, but child-lifecycle management (watchdog propagation, convergence checking) scales with both.

The $d_{\max}$ parameter mitigates unbounded depth but does not bound branching. A per-node spawn budget or a system-wide spawn rate limiter could complement depth limits, but these mechanisms introduce additional complexity and have not yet been formalized under RSP. Future work should evaluate bounded branching factor models in which each node's maximum child count is a function of its depth and safety envelope, analogous to the per-authority depth budget $\Delta(p)$ defined in Section~\ref{sec:rsp-bounds-engine}.

\subsection{Compromised Purpose Layer}
\label{sec:limitations-purpose-compromise}

RSP's accountability guarantees rest on the integrity of the Purpose Layer validator. An agent that can successfully falsify its spawn justification to the Purpose Layer---presenting a valid-appearing distinct-capability claim for a child that is functionally redundant with its own capabilities---can route unauthorized spawn events through the protocol's first gate. Once a child is spawned with a falsified Purpose declaration, the child is structurally indistinguishable from a legitimately spawned agent: it carries a valid immutable parent pointer and a signed warrant from the Fact Layer.

This limitation is not unique to RSP; it is a manifestation of a broader problem in accountable AI: the system cannot independently verify the semantic accuracy of a natural language capability justification. The Purpose Layer validator performs structural checks (redundancy detection against the task registry, context budget analysis) but does not perform deep semantic verification of whether the declared child capability is genuinely necessary or uniquely provided. A sophisticated adversarial agent could, in principle, construct spawn justifications that pass structural validation while routing functionally unnecessary children into the spawn tree.

The drift detection mechanism described in Section~\ref{sec:rsp-fact-layer} provides a secondary defense: even a child spawned with a falsified Purpose is subject to behavioral monitoring against its authorized scope $\sigma_c$. If the child's observed behavior diverges from $\sigma_c$, the Bailout Protocol is triggered. The Purpose Layer compromise is therefore not an unmitigated vulnerability; it is a vulnerability with a compensating control. However, the control depends on the accuracy and observability of the Observation Layer, which itself may be compromised. Defending against coordinated compromise of multiple BX3 layers is an open problem.

Formal methods for the verification of Purpose Layer integrity under adversarial conditions---particularly in the case where the parent agent itself is the source of the falsified justification---remain an active area of future research. The RSP provides structural accountability; it does not eliminate the possibility of strategic abuse by sophisticated actors.

\subsection{Future Work}
\label{sec:future-work}

Four directions for future work follow directly from the limitations above.

\textbf{Formal verification.} RSP's guarantees---finite termination, constant-time escalation, immutable parent pointer integrity---should be formalized in a proof assistant (e.g., Coq or Isabelle/HOL) to establish machine-checked proofs of correctness independent of runtime implementation. Prior work on supervisory control for hierarchical multi-agent systems provides a relevant formal foundation.\cite{ouhamma2024hmas}

\textbf{Adversarial robustness.} Extension of RSP to settings where parent agents are actively adversarial requires integration of cryptographic attestation into the spawn justification pathway, such that the Purpose Layer validator can verify not only the structural correctness of the spawn proposal but also the integrity of the parent agent's own state at spawn time. This is architecturally related to the challenge of Observation Layer integrity described above and may benefit from shared solutions.

\textbf{Bounded branching factor models.} Formal specification and evaluation of per-node spawn budgets as a complement to depth limits. This direction is particularly relevant for high-density IoT deployments (e.g., SLV agricultural sensor networks) where the branching factor is determined by physical sensor topology rather than agentic reasoning decisions.

\textbf{Domain-specific calibration.} RSP's parameters ($d_{\max}$, spawn timeouts, watchdog intervals, convergence definitions) require calibration to specific deployment contexts. Agricultural deployments managing real-time irrigation decisions have different latency tolerances and safety requirements than corporate deployments managing document workflows. A calibration methodology that relates RSP parameters to domain-specific risk profiles would extend the protocol's practical applicability.

\end{document}
